\chapter{Beyond linear and stationary agents}\label{ch:beyondlinear}

Chapters~\ref{ch:p1} and~\ref{ch:p2} work with agents whose rules are linear filters with fixed memory and fixed weight. Real agents cap positions, trade only beyond thresholds, adapt their gains, and keep learning while the operator watches. This chapter shows that the safety side of the thesis survives all of this under one condition: a bound on the agents' total sensitivity. It then asks what remains open when agents respond strategically to the guard itself.

\section{A robust safety bound}\label{sec:bl_robust}

Return to the economy of Definition~\ref{def:cohort} and let everything about the agents vary over time.

\begin{assumption}[Bounded total sensitivity]\label{ass:sector}
For each cohort $i$ and period $t$, the memory rate $r_i(t)\in[0,1)$ is arbitrary, the weight satisfies $|q_i(t)|\le\bar q_i$, and the response rule satisfies the sector condition $|\varphi_{i,t}(y)|\le s_i|y|$ for all $y$. Write $\bar\Phi=\sum_i\bar q_i s_i$ for the total sensitivity.
\end{assumption}

The sector condition covers saturating momentum ($s_i=1$ for $\tanh$), position limits, threshold and dead-zone rules ($y\mathbf 1\{|y|>\tau\}$), and any rule whose response never exceeds a multiple of its signal. It does not require the rule to be smooth, monotone, stationary or known.

\begin{theorem}[Robust safety bound]\label{thm:sector}
Under Assumption~\ref{ass:sector}, if $|a|+\kappa\bar\Phi<1$ and the economy starts at rest, then for every shock sequence with $|w_t|\le\bar w$,
\[
\sup_{t\ge0}|x_t|\;\le\;\bar w\,\frac{\kappa}{1-|a|-\kappa\bar\Phi}.
\]
When $a\ge0$, this is the worst-case loss bound of the linear stationary economy with weights $\bar q_is_i$, so the summed-gain guard of Chapter~\ref{ch:p2} computed at total sensitivity $\bar\Phi$ is safe for every economy satisfying the assumption. \hfill\Proved
\end{theorem}

\begin{proof}
Let $Z_t=\max_{s\le t}|x_s|$. Each memory state is a convex combination of its previous value and the current state, so $|y_{i,t}|\le r_i(t)|y_{i,t-1}|+(1-r_i(t))|x_t|\le\max(|y_{i,t-1}|,|x_t|)$, and by induction from $y_{i,0}=0$, $|y_{i,t}|\le Z_t$. Then
\[
|x_{t+1}|\le|a|\,|x_t|+\kappa\Bigl(\sum_i|q_i(t)|\,|\varphi_{i,t}(y_{i,t})|+|w_{t+1}|\Bigr)\le(|a|+\kappa\bar\Phi)\,Z_t+\kappa\bar w .
\]
Let $Z^*=\kappa\bar w/(1-|a|-\kappa\bar\Phi)$. If $Z_t\le Z^*$ then $|x_{t+1}|\le(|a|+\kappa\bar\Phi)Z^*+\kappa\bar w=Z^*$, so $Z_{t+1}\le Z^*$; and $Z_0=0$. For the second claim, with $a\ge0$ the linear stationary economy with weights $\bar q_is_i$ is a positive system, its impulse response is non-negative, and its worst-case peak response to shocks bounded by $\bar w$ is $\bar w\sum_th_t=\bar w\,\kappa/(1-a-\kappa\bar\Phi)$ (Chapter~\ref{ch:framework}, equation~\eqref{eq:fw_S}). The bound of the theorem therefore coincides with the quantity the guard controls (Lemma~\ref{p2:thm:guard}).
\end{proof}

\paragraph{What the theorem says.} An operator who wants a safe admission level does not need to know how many agents there are, what their memories are, whether their rules are linear, or whether they change over time. It needs one number, an upper bound on the agents' total sensitivity $\bar\Phi$. The bound is tight: when the agents are linear, stationary, positive and at their maximal weights, it is attained. Nonlinearity, memory changes and negative feedback can only make the market safer than the bound says.

\begin{remark}[What the theorem does not cover]
Rules that amplify their signal, $|\varphi(y)|>s|y|$ for the declared $s$, break the bound; so do changes in the market's own persistence $a$ or impact $\kappa$, which the operator must certify separately (Chapter~\ref{ch:p2}, Section~\ref{p2:sec:infosafety}). The bound controls the peak of the price path, which is what a fixed position's loss depends on; dynamic positions are handled by the balance-sheet guard of Chapter~\ref{ch:p2}.
\end{remark}

\begin{proposition}[Several assets]\label{prop:multiasset}
Let $x_t\in\R^d$ evolve as $x_{t+1}=Ax_t+\mathcal K\bigl(\sum_iQ_i(t)\varphi_{i,t}(y_{i,t})+w_{t+1}\bigr)$, with memory states $y_{i,t}=R_i(t)y_{i,t-1}+(I-R_i(t))x_t$, $R_i(t)$ diagonal with entries in $[0,1)$, sector conditions $|\varphi_{i,t}(y)|\le s_i|y|$ componentwise, and $\|Q_i(t)\|_\infty\le\bar q_i$. If $\|A\|_\infty+\|\mathcal K\|_\infty\sum_i\bar q_is_i<1$, then $\sup_t\|x_t\|_\infty\le\bar w\,\|\mathcal K\|_\infty/(1-\|A\|_\infty-\|\mathcal K\|_\infty\sum_i\bar q_is_i)$ for $\|w_t\|_\infty\le\bar w$. \hfill\Proved
\end{proposition}

\begin{proof}
The proof of Theorem~\ref{thm:sector} goes through with absolute values replaced by the sup norm and scalars by induced sup norms: each component of $y_{i,t}$ is a convex combination of its past value and the corresponding component of $x_t$.
\end{proof}

The multi-asset bound is conservative, because induced norms ignore cancellation across assets; a sharp version for positive cross-impact structures, matching the leader--satellite analysis of Chapter~\ref{ch:p2}, is a \Target{}.

\section{Certificates expire: learning agents}\label{sec:bl_learning}

Theorem~\ref{thm:sector} allows weights to change, provided they stay below their bounds. When agents learn, the bounds themselves are the issue: a certificate issued today says nothing about weights next month. The theorem turns this into a computable lifetime.

\begin{corollary}[Certificate lifetime]\label{thm:lifetime}
Suppose a certificate at time $0$ establishes $|q_i(0)|\le q_i^+$ for each cohort, and the guard's loss budget is met for every total sensitivity up to $\Phi_{\max}$, that is $\kappa\bar w/(1-a-\kappa\Phi_{\max})$ times the position equals the budget. If each weight changes by at most $\gamma_i$ per period, $|q_i(t+1)-q_i(t)|\le\gamma_i$, then the guard remains safe for at least
\[
T_{\rm valid}=\Bigl\lfloor\frac{\Phi_{\max}-\sum_iq_i^+s_i}{\sum_i\gamma_is_i}\Bigr\rfloor
\]
periods, whatever the agents' memories and rules within the sector bounds. \hfill\Proved
\end{corollary}

\begin{proof}
On $[0,T_{\rm valid}]$, $|q_i(t)|\le q_i^++\gamma_iT_{\rm valid}=:\bar q_i$ and $\sum_i\bar q_is_i\le\Phi_{\max}$. Theorem~\ref{thm:sector} applies on that horizon, since its induction uses only periods up to the current one.
\end{proof}

The lifetime is the certificate's safety margin divided by the agents' speed of learning. It cannot be lengthened in general: if all weights rise at their maximal rates and then stay at the resulting level, the economy after $T_{\rm valid}$ periods is the linear economy at $\Phi_{\max}$, whose worst-case loss under persistent shocks approaches the budget. Two policy consequences follow. Recertification frequency should be set by how fast the agents in a market adapt, not by the calendar; and a market whose agents learn faster needs either more frequent certificates or a larger safety margin, a trade-off the formula prices.

\paragraph{Identification under learning, by contrast.} Chapter~\ref{ch:p1}, Section~\ref{p1:sec:learners}, shows that trained online learners obey the same $\delta^{-4}$ law for separating memory rates when the competing populations' weights match, and that populations whose learned weights differ are told apart by those weights instead. Learning therefore leaves identification as hard as before while making certification only perishable, not harder per renewal.

\section{Agents who respond to the guard}\label{sec:bl_strategic}

A guard changes the environment it certifies. If agents learn that the operator throttles admission when measured sensitivity is high, they may change how they trade.

\begin{conjecture}[Guard-consistent equilibrium]\label{conj:strategic}
Suppose each cohort chooses its weight $q_i\in[0,\bar q_i]$ to maximize expected trading profit net of a convex cost, given an admission rule that throttles the dealer's position when certified total sensitivity exceeds $\Phi_{\max}$, and given the equilibrium reversion supplied by dealers (Proposition~\ref{p2:prop:equil}). Then (i) an equilibrium exists; (ii) in every equilibrium the guard is safe; and (iii) equilibrium total sensitivity is weakly lower than without the guard. \hfill\Conj
\end{conjecture}

The conjecture is plausible because the guard makes high sensitivity costly to the agents that create it, through lower admitted liquidity. It can fail if agents can hide sensitivity from the certificate, for example by trading only in states the reset experiments do not visit; this is the reason the conjecture is stated with a certificate of the total sensitivity rather than of the summed response at the current state. A proof would combine Theorem~\ref{thm:sector}, which makes safety depend only on $\bar\Phi$, with a fixed-point argument on the agents' weights.
